\section{Consequences for the spectral and geometric programmes}
This section separates consequences of the zero measure from statements
about an individual eigenvector, a cohomology group, or an automorphic
representation. Those objects require different comparisons.

\subsection{Zeta cycles and the lowest vector}
The characterization of zeta cycles in
\citet[Theorem~1.1, p.~4, and Theorem~6.4, pp.~52--53]{cc2021}
concerns the zeros on the critical line. The verified Riemann hypothesis
makes that set the complete set of nontrivial zeta zeros. This deduction
does not give their simplicity or recover their multiplicities from
membership in the set of cycles. Likewise, the spectral counting
asymptotic of \citet[Theorem~6.1, p.~13]{cm2021} remains a counting
asymptotic; identifying its individual eigenvalues requires a separate
operator comparison.

\citet[Theorem~5.10, p.~23]{ccm2025} prove a more specific finite
statement. If the lowest eigenvalue of their finite Weil matrix is
simple and its eigenvector $\xi_N$ is even, then the rank-one
perturbation of the scaling generator is self-adjoint for the
quotient metric, and its regularized determinant is
$-i\lambda^{-iz}\widehat\xi_N(z)$ in their normalization. In
particular, the zeros of that transform are real. The theorem already
proves this implication. Neither the Riemann hypothesis nor a trace
limit removes its hypotheses about the lowest eigenvector. The same
paper explicitly lists eventual simplicity, evenness, and approximation
of that vector among its missing steps (p.~32).

Here is the exact elementary estimate needed for the first of those
steps. For $L>0$, let $\mathcal A_L$ be a nonnegative self-adjoint
operator with compact resolvent on $L^2(-L,L)$, commuting with
reflection and complex conjugation. Write its eigenvalues, with
multiplicity, as $a_1(L)\leq a_2(L)\leq\cdots$, and let $q_L$ be its
closed quadratic form. These hypotheses apply to the localized Weil
form after the Riemann hypothesis. All functions are extended by zero.

\begin{proposition}[The ground-vector estimate]
\label{prop:ground-vector-ratio}
Let $v_L$ be a real even unit vector in the form domain, and suppose
that $a_2(L)>0$. If $r_L=q_L(v_L)/a_2(L)<1$, the lowest eigenspace
is simple and even. Its real unit vector $u_L$, chosen with
$\langle u_L,v_L\rangle>0$, satisfies
\[
 1-|\langle u_L,v_L\rangle|^2\leq r_L,
 \qquad \|u_L-v_L\|_2^2\leq2r_L.
\]
Consequently, if $v_L\to K/\|K\|_2$ and $r_L\to0$, then
$u_L\to K/\|K\|_2$ in $L^2(\R)$.
\end{proposition}
\begin{proof}
The minimum principle gives $a_1(L)\leq q_L(v_L)<a_2(L)$, so the
lowest eigenvalue is simple. If its eigenvector were odd, the even
vector $v_L$ would be orthogonal to it and would have quotient at least
$a_2(L)$, a contradiction. Expansion in an orthonormal eigenbasis
gives
$q_L(v_L)\geq a_2(L)(1-|\langle u_L,v_L\rangle|^2)$.
The phase can be chosen real and positive. For $0\leq b\leq1$,
$2(1-b)\leq2(1-b^2)$, proving the norm bound. The last assertion
follows by the triangle inequality.
\end{proof}

For the Riemann kernel $K$, normalized by
$\widehat K(z)=\Xi(z)=\xi(1/2+iz)$, the prolate trial construction
has the estimate
\[
 q_{\log\lambda}(v_\lambda)
       \ll_\varepsilon\lambda^{9+\varepsilon}e^{-4\pi\lambda^2}.
\]
Thus a sufficient, still open, lower bound is
\begin{equation}\label{eq:ground-second-eigenvalue-gate}
 a_2(\log\lambda)\geq c\lambda^{9+\alpha}e^{-4\pi\lambda^2}
 \quad\text{for some }c,\alpha>0\text{ and all sufficiently large }\lambda.
\end{equation}
This is a lower bound on the second eigenvalue of the Weil operator,
not a bound on the concentration eigenvalues of the raw prolate
operator. The trace formula proved above supplies neither
\eqref{eq:ground-second-eigenvalue-gate} nor the weaker condition
$r_\lambda\to0$. The known odd trial vector bounds
$a_2(\log\lambda)\ll_\varepsilon
\lambda^{13+\varepsilon}e^{-4\pi\lambda^2}$, so a power lower bound
of the displayed form must have $\alpha\leq4$. These trial estimates
are the propositions ``Leakage controls the original trial'' and
``An odd prolate competitor'' in
\citet[section ``The prolate trials and their Rayleigh upper bounds'']
{simplicitynotes}.

There is a further distinction between $L^2$ convergence and entire
convergence of Fourier transforms on expanding intervals. The former
does not imply the latter by the Cauchy--Schwarz inequality: evaluation
at a nonreal point incurs a factor growing exponentially with $L$.
For the finite Weil determinants one instead has real-rooted polynomial
approximants. The convergence theorem of
\citet[Corollary~1.3, p.~110]{cluniekuijlaars1994}, applied after
a diagonal polynomial approximation, upgrades convergence in measure
on a real interval where $\Xi\ne0$ to local uniform entire convergence.
This route still requires $r_\lambda\to0$ and the finite determinant
hypotheses. Write $a_e(L)$ and $a_o(L)$ for the even and odd minima.
The proposition ``A subsequence criterion for simplicity'' in
\citet[section ``The lowest Weil eigenvector and the simplicity question'']
{simplicitynotes} shows that even a subsequence with
$r_\lambda\to0$ would force simplicity of all zeta zeros: along
that sequence $a_e/a_o\leq r_\lambda$, so the fixed relative odd
gap identifies the limiting normalized transform with the product
over distinct zeros, whereas the $L^2$ limit identifies it with
$\Xi/\Xi(0)$. The two products can agree only when every
multiplicity is one. The trace limit alone proves no such ratio assertion.

A different sufficient condition uses the even and odd minima
$a_e(L)$ and $a_o(L)$. Its quantitative hypothesis is
\begin{equation}\label{eq:relative-odd-gap}
 a_e(L)<a_o(L),\qquad
 1-a_e(L)/a_o(L)\geq\kappa>0
 \quad\text{eventually}.
\end{equation}
In the Weil Hilbert space, write the compressed primitive and the
compact inclusion operator in parity blocks as
\[
 C_L=\begin{pmatrix}0&X_L^*\\X_L&0\end{pmatrix},\qquad
 T_L=\begin{pmatrix}T_{L,+}&0\\0&T_{L,-}\end{pmatrix},\qquad
 S_\pm=I-a_e(L)T_{L,\pm}.
\]
The quotient primitive is represented by
\[
 B_L=\begin{pmatrix}
 0&S_+^{1/2}X_L^*S_-^{-1/2}\\
 S_-^{-1/2}X_LS_+^{1/2}&0
 \end{pmatrix}.
\]
The fixed lower bound on $S_-$ in
\eqref{eq:relative-odd-gap} is the estimate that controls this
comparison. Its known consequence is convergence to the product
$\prod_{\gamma>0,\ \mathrm{distinct}}(1-z^2/\gamma^2)$.
This is the proposition ``The parity comparison'' in the same
section of \citet{simplicitynotes}.
Identification with $\Xi(z)/\Xi(0)$ additionally requires simplicity.
An abstract compact example satisfying the rank-one commutator
identity but with diverging $\|B_L\|$ shows why the commutator alone
does not supply this estimate
\citep[proposition ``Exact commutator without uniform comparison'']
{commutatornotes}. It is not a counterexample for the
Weil operator. None of these conditional assertions is changed by
the scalar trace limit.

\subsection{Effective cohomology and a positive Hilbert defect}
There is a precise elementary part of the analogy with a finite-field
trace formula. An effective cohomological remainder is the trace on
an honest vector space, with positive integer algebraic multiplicities.
It need not be a positive functional on a star algebra.

\begin{proposition}[Power traces and effectivity]
\label{prop:effective-power-traces}
Let $F$ be an endomorphism of a finite-dimensional complex vector
space. If $|\Tr F^n|\leq CB^n$ for every integer $n\geq1$, with
$B>0$, then every eigenvalue of $F$ has absolute value at most $B$.
This conclusion does not imply that $F$ is semisimple or that it
admits a faithful positive realization with a prescribed adjoint.
\end{proposition}
\begin{proof}
For sufficiently small $|z|$, Jordan normal form gives
\[
 \sum_{n\geq1}\Tr(F^n)z^n
     =\sum_{\lambda\ne0}m_\lambda\frac{\lambda z}{1-\lambda z}.
\]
The coefficient bound makes the left side holomorphic on
$|z|<B^{-1}$. If $|\lambda|>B$, the right side has a pole at
$z=\lambda^{-1}$; its multiplicity $m_\lambda$ is a positive integer,
so that pole cannot cancel with a different eigenvalue. This is a
contradiction. For the second assertion, take a nonzero nilpotent
$3\times3$ Jordan matrix $N$ and $F=c\exp N$, $c\ne0$.
Then $\Tr(F^n)=3c^n$ although $F$ is not semisimple. In the algebra
$\C[N]$, prescribe $N^*=N$. The ordinary trace gives
$\Tr(a^*a)=3|a_0|^2$ for $a=a_0+a_1N+a_2N^2$, so this positive
functional is not faithful: $\Tr(N^*N)=0$ although $N\ne0$.
More generally, a Hilbert-space star representation sends $N$ to
a self-adjoint nilpotent operator, which is zero by the spectral
theorem. Such a representation cannot be faithful.
\end{proof}

Consider a sign-only purity argument in which the cohomological
remainder is realized, after a fixed global sign, on a Gysin kernel
and a tensor power of compactly supported cohomology. This realization
is the hypothesis of the comparison here. Its effective class makes
the preceding power-sum argument possible. A virtual difference of
two representations could cancel equal eigenvalues, and would not
justify that argument without further information. By comparison,
$S_L-S_LQ_\Lambda S_L$ is a positive operator because it is
$S_L(1-Q_\Lambda)S_L$; its positivity is an order statement in a
Hilbert space. Effectivity of a Gysin kernel does not construct this
Hilbert structure, and the analytic defect does not construct a
Gysin map over $\operatorname{Spec}\Z$.

The distinction is consistent with the precise questions in
SGA~6, Expos\'e~XIV, \S6.1--6.4
\citep[pp.~682--685]{sga6}. There the cohomological Riemann--Roch
diagram is conditional on defining the Gysin homomorphism; the text
states the available regularity and purity hypotheses and asks about
the remaining cases. It does not assert a universal Gysin construction
for singular schemes. Transport from a finite-field trace formula
therefore needs both a defined characteristic-zero realization and
the quantitative comparison of its finite sections. The two end
boundaries and the inverse Gram estimate above perform the latter
analytic task; they are not a substitute for the former geometric one.

\subsection{The abelian range of Frobenius lifts}
For a smooth proper scheme over $\Z[1/N]$ equipped with commuting
Frobenius lifts, Borger's theorem gives abelian Galois action and a
sum of Tate characters $\Q_\ell(-m)$ after a cyclotomic base extension
\citep[Theorem~0.1 and Corollary~6.2]{borger2009}. Purity in degree
$i$ forces every Tate exponent $m$ to satisfy $2m=i$: evaluate a
good geometric Frobenius that splits in that cyclotomic field, whose eigenvalue
has both absolute value $p^m$ and weight $p^{i/2}$. Thus odd
cohomology vanishes, and the Tate-normalized even cohomology has
finite abelian image. Direct sums, tensor products, duals, and
subquotients preserve that property.

Conversely, a finite abelian rational $\ell$-adic representation
factors through a cyclotomic Galois group. The finite \hbox{\'etale}
scheme $\operatorname{Spec}\Z[\zeta_N,1/N]$ has commuting good-prime
Frobenius lifts and its degree-zero cohomology is the regular
representation. Maschke's theorem realizes the given representation
as a summand of finitely many copies. Hence this generated category
is exactly the category of finite abelian image representations.
The Dirichlet trace formula applies component by component after
Tate normalization, with the original twists restored by translating
the Mellin variable. This proves no analogous statement for a
nonabelian representation merely by realizing it in degree-zero
cohomology of a finite \hbox{\'etale} cover.

For example, the standard two-dimensional representation of $S_3$
does not kill its commutator subgroup $A_3$ and is outside this
abelian category. Over $K=\Q(\sqrt{-3})$, however, the extension
$K(\sqrt[3]{269})/K$ is cyclic. At a good prime ideal $\mathfrak p$
of norm $q$, its unique relative Frobenius lift sends
$\sqrt[3]{269}$ to $\zeta_\mathfrak p\sqrt[3]{269}$, where
$\zeta_\mathfrak p\equiv269^{(q-1)/3}\pmod{\mathfrak p}$.
These lifts commute because the deck group is cyclic. This identifies
the relative Hecke character; it does not place its zeros on the line.

\subsection{A Godement--Jacquet finite-section question}
The theta construction needed beyond $\mathrm{GL}_1$ is already
available. Fix an irreducible cuspidal automorphic representation
$\pi$ of $\mathrm{GL}_n(\mathbb A_\Q)$ and a place $v$.
Put $F=\Q_v$, let $\beta$ be a matrix coefficient of $\pi_v$, and
let $\Phi\in\mathcal S(\mathrm M_n(F))$. The local
Godement--Jacquet integral is
\[
 Z(s,\Phi,\beta)=\int_{\mathrm{GL}_n(F)}\Phi(g)\beta(g)
                  |\det g|^{s+(n-1)/2}\,dg.
\]
Its determinant-fibre pushforward is
\[
 \phi_{\Phi,\beta}(x)=|x|^{n/2}
                  \int_{\det g=x}\Phi(g)\beta(g)\,d_xg.
\]
With compatible quotient measures,
$\int\phi_{\Phi,\beta}(x)|x|^{s-1/2}d^\times x=Z(s,\Phi,\beta)$.
These formulas and their local convergence are given by
\citet[Proposition~3.2 and Definition~3.3, pp.~16--17]{jiangluo2023};
archimedean test vectors attaining the local $L$-factor are treated
by \citet[Theorem~4.22]{humphries2025}.

On the associated global $\pi$-Schwartz space, Jiang and Luo prove
the Poisson formula
\[
 \Theta_\pi(x,\phi)=\sum_{\alpha\in\Q^\times}\phi(\alpha x)
     =\Theta_{\widetilde\pi}(x^{-1},\mathcal F_\pi\phi)
\]
for cuspidal $\pi$ \citep[Theorem~4.7, p.~27]{jiangluo2023}.
Their critical-line quotient theorem has multiplicity limited by
the weight of the chosen Hilbert space
\citep[Theorem~8.2, p.~50]{jiangluo2023}. It is not the all-zero
finite-cutoff formula considered here. Jiang and Li's kernel
characterization of finite-part $L$-zeros is another established
comparison \citep[Theorems~6.9 and 6.13, pp.~34 and 39]{jiangli2024}.

The following finite-section construction makes positivity exact. The
zero-measure limit remains a separate question, whose estimates are
stated after the construction.

\subsection{A positive windowed Godement--Jacquet cutoff}
\label{glnw:section}

There is a concrete positive finite section on the cuspidal automorphic
quotient, provided the projection is formed after the norm window is
applied. This construction has a different normalization from the
unwindowed theta-image projection.

Let $G=\mathrm{GL}_n$, with $n\geq2$, and fix an irreducible
unitary cuspidal representation $\pi$ of $G(\mathbb A_\mathbb Q)$. Use the
Godement--Jacquet determinant pushforward and the $\pi$-theta map
described above. Write $C_\mathbb Q^1$ for the compact norm-one
idele class group, choose a section $x(t)$ of the idele norm with
$|x(t)|=e^t$. Choose compatible quotient measures, with
$C_\mathbb Q^1$ of volume one. The untwisted
norm channel is
\[
 E_\pi\phi(t)=\int_{C_\mathbb Q^1}
                   \Theta_\pi(x(t)u,\phi)\,du .
\]
The rapid decay in both norm directions in
\citet[Proposition~8.1, p.~49]{jiangluo2023} puts this function
in $L^2(\mathbb R,dt)$. Define the global zeta integral, initially
in its half-plane of absolute convergence, by
\[
 Z(s,\phi)=\int_{\mathbb A_\mathbb Q^\times}
               \phi(x)|x|^{s-1/2}\,d^\times x.
\]
For a determinant pushforward this is the global
Godement--Jacquet integral, with the fixed matrix coefficient and
finite-place data understood. Unfolding gives
\begin{equation}\label{glnw:mellin}
 \int_{\mathbb R}E_\pi\phi(t)e^{(s-1/2)t}\,dt=Z(s,\phi).
\end{equation}
The rapid norm decay and the global continuation extend this equality
to every $s$. The local zeta-integral ideal of
\citet[Theorem~2.4(2), p.~10]{jiangluo2023} supplies the usual
$L(s,\pi)$ factor. In particular, wherever the complementary factor is
holomorphic, a zero of $L(s,\pi)$ gives the corresponding full Mellin
jets of $E_\pi\phi$. These are statements about the unwindowed theta
map, not yet about a finite-section limit.

Here is one explicit choice of finite input spaces. On the real matrix
space $\mathrm M_n(\mathbb R)$, of dimension $m=n^2$, use the unitary
Fourier transform with kernel $(2\pi)^{-m/2}e^{-i\operatorname{tr}(X^T\Xi)}$.
Let $T_\Lambda=B_\Lambda=\{X:\|X\|_{\rm F}<\Lambda\}$,
where $\|X\|_{\rm F}^2=\operatorname{tr}(X^TX)$ is the Frobenius norm. The ordinary concentration operator
\[
 A_\Lambda^{\rm mat}
  =1_{T_\Lambda}\mathcal F^{-1}1_{B_\Lambda}
                         \mathcal F1_{T_\Lambda}
\]
is positive and trace class, with
\[
 \Tr A_\Lambda^{\rm mat}
       =(2\pi)^{-m}|T_\Lambda||B_\Lambda|.
\]
For a fixed $0<\theta<1$, its spectral selection on $[\theta,1]$
has finite rank, at most $\theta^{-1}\Tr A_\Lambda^{\rm mat}$.
Apply to these vectors a composition
$M_{\eta_T}\mathcal F^{-1}M_{\eta_B}\mathcal F$, where the two
multipliers are smooth and compactly supported inside their
respective balls. The resulting vectors are in the matrix Schwartz
space. This composition is used to enter that space; it asserts no
simultaneous exact time and frequency support. Fix the finite-place
Schwartz data and a finite set of matrix coefficients, form the
determinant pushforwards, and then form their theta images. A fixed
finite set of compact-group types may also be imposed on the matrix
selection. The construction remains finite dimensional.

Let $Y_\Lambda:\mathbb C^{N_\Lambda}\longrightarrow
\mathcal H_d=L^2(\mathbb R)\otimes\mathbb C^d$ be the synthesis map
of any such finite set of theta columns. The integer $d$ counts the
chosen coefficient channels; it is not the degree $n$ of the
standard $L$-function. Put
\[
 S_L=1_{[-L,L]}\otimes I_d,\qquad
 \widetilde G_{\Lambda,L}=Y_\Lambda^*S_LY_\Lambda,\qquad
 \widetilde Q_{\Lambda,L}
 =S_LY_\Lambda\widetilde G_{\Lambda,L}^{\dagger}
                         Y_\Lambda^*S_L ,
\]
where $\dagger$ denotes the Moore--Penrose inverse. Equivalently,
remove dependent windowed columns and invert their full Gram
matrix. The scale $L=L(\Lambda)$ can be any prescribed positive
function, for example $L=\log\Lambda$; no comparison between matrix
concentration and norm concentration is needed for the next assertion.

\begin{proposition}[Exact finite-section positivity]
\label{glnw:positive}
The operator $\widetilde Q_{\Lambda,L}$ is the orthogonal projection
onto $\operatorname{Ran}(S_LY_\Lambda)$ and satisfies
$\widetilde Q_{\Lambda,L}\leq S_L$. Hence
$\widetilde D_{\Lambda,L}=S_L-\widetilde Q_{\Lambda,L}$ is an
orthogonal projection. For $a\in C_c^\infty(\mathbb R)$, set
$a^*(t)=\overline{a(-t)}$ and $h=a*a^*$. Then
\begin{equation}\label{glnw:positive-trace}
 2Ld\,h(0)
 -\Tr\bigl(\widetilde Q_{\Lambda,L}(V(h)\otimes I_d)\bigr)
 =\bigl\|\widetilde D_{\Lambda,L}
                   (V(a)\otimes I_d)\bigr\|_{\mathrm{HS}}^2
 \geq0.
\end{equation}
In particular this defect is positive at each finite scale.
\end{proposition}

\begin{proof}
For a finite-dimensional synthesis map $B$, the operator
$B(B^*B)^\dagger B^*$ is the orthogonal projection onto its range;
this follows directly from a singular-value decomposition.
Apply this to $B=S_LY_\Lambda$. Its range lies in that of $S_L$,
so the two projections are nested and their difference is a
projection. The kernel of $S_L(V(a)\otimes I_d)$ gives
\[
 \|S_L(V(a)\otimes I_d)\|_{\mathrm{HS}}^2
        =2Ld\,\|a\|_2^2=2Ld\,h(0).
\]
Since $S_L=\widetilde Q_{\Lambda,L}+\widetilde D_{\Lambda,L}$
is an orthogonal sum, the corresponding Hilbert--Schmidt norms
add in squares. The first summand is
$\Tr(\widetilde Q_{\Lambda,L}(V(h)\otimes I_d))$ because
$V(h)=V(a)V(a)^*$. This proves the formula and justifies all
traces without assigning an infinite trace to an unlocalized
convolution operator.
\end{proof}

\begin{proposition}[Realization on the cuspidal quotient]
\label{glnw:quotient}
The construction of Proposition~\ref{glnw:positive} is realized
isometrically on the cuspidal subspace of
$G(\mathbb Q)\backslash G(\mathbb A_\mathbb Q)$ generated by the
unitary norm twists of $\pi$, with a fixed finite set of
cuspidal coefficient channels. It does not lie in the single
fixed-central-character representation $\pi$ alone.
\end{proposition}

\begin{proof}
Let
\[
 G^1=\{g:|\det g|_{\mathbb A}=1\},\qquad
 X^1=G(\mathbb Q)\backslash G^1 .
\]
For $t\in\mathbb R$, let $a(t)$ have real component
$e^{t/n}I_n$ and finite components $I_n$. It is central and
$|\det a(t)|=e^t$. Each $g\in G(\mathbb A_\mathbb Q)$ has the unique decomposition
$g=a(t)g_1$ with $t=\log|\det g|$ and $g_1\in G^1$.
The product formula puts $G(\mathbb Q)$ inside $G^1$, so this
gives a product decomposition of the automorphic quotient.
Normalize Haar measures so that its measure is $dt\,dx_1$.

Choose orthonormal cuspidal vectors
$e_1,\ldots,e_d\in L^2(X^1)$ from the restriction of the
cuspidal realization of $\pi$. Define
\[
 (\mathcal U f)(a(t)g_1)=\sum_{j=1}^d f_j(t)e_j(g_1).
\]
Fubini and orthonormality show that $\mathcal U$ is an isometry
from $\mathcal H_d$. Every proper-parabolic unipotent subgroup
lies in $G^1$ and commutes with the norm section in the required
constant-term calculation. The vanishing of the constant term
of each $e_j$ therefore proves that $\mathcal U f$ is cuspidal.
The identity extends from smooth tensors to their $L^2$ closure.

For the right regular representation $R$, one has
\[
 R(a(-s))\mathcal U f=\mathcal U(f(\,\cdot-s)).
\]
Thus $\mathcal U(V(h)\otimes I_d)\mathcal U^*$ is the integrated
norm action on the range of $\mathcal U$. Transport
$S_L$ and $\widetilde Q_{\Lambda,L}$ by this isometry. The first
becomes the norm-window projection with the additional fixed
projection onto $\operatorname{span}\{e_1,\ldots,e_d\}$; the
second is finite rank. Formula~\eqref{glnw:positive-trace} is
unchanged.

Fourier transformation in $t$ decomposes the norm translations
into unitary characters of $\mathbb R$. Together with the
fixed $G^1$ realization, these are the unitary twists of $\pi$
by powers of $|\det|$. A fixed irreducible representation
with a fixed central character has scalar action of $a(t)$.
It cannot itself contain this regular norm action. The
finite-channel space is also not asserted to be invariant
under all of $G$; only the norm action is used in the trace.
\end{proof}

\begin{proposition}[Comparison with the unwindowed projection]
\label{glnw:comparison}
Let $Q_\Lambda$ be the orthogonal projection onto
$\operatorname{Ran}Y_\Lambda$, and suppress the indices.
Then
\[
 \widetilde Q=\mathbf1_{(0,\infty)}(SQS),\qquad
 0\leq SQS\leq\widetilde Q\leq S .
\]
Moreover
\begin{equation}\label{glnw:normalization-loss}
 \Tr(\widetilde Q-SQS)
  =\operatorname{rank}(SQ)-\Tr(QS)
  \leq\|(1-S)Q\|_{\mathrm{HS}}^2 .
\end{equation}
Equality holds in the last inequality when $S$ is injective
on $\operatorname{Ran}Q$.
\end{proposition}

\begin{proof}
The positive finite-rank operator $SQS=(SQ)(SQ)^*$ has range
$\operatorname{Ran}(SQ)=\operatorname{Ran}(SY)$. Its support
projection is therefore $\widetilde Q$, and its eigenvalues
on that range belong to $(0,1]$. This proves the order
identities. Taking traces gives the equality in
\eqref{glnw:normalization-loss}. Since
$\operatorname{rank}(SQ)\leq\operatorname{rank}Q$,
the remaining inequality follows from
$\operatorname{rank}Q-\Tr(QS)=\Tr(Q(1-S)Q)$.
\end{proof}

\begin{proposition}[Exact normalization criterion]
\label{glnw:exact-normalization}
Put $T=SQS$, $\Delta=\widetilde Q-T$, and let $W$ be bounded.
The Moore--Penrose inverse is zero on $\ker T$. Then
\begin{equation}\label{glnw:inverse-boundary-criterion}
 \Delta=T^{\dagger/2}(T-T^2)T^{\dagger/2},\qquad
 \Tr(\Delta WW^*)
   =\|(1-S)QS T^{\dagger/2}W\|_{\mathrm{HS}}^2 .
\end{equation}
If $T\geq\kappa\widetilde Q$ for some $\kappa>0$, then
\begin{equation}\label{glnw:window-gap-bound}
 \Tr(\Delta WW^*)
 \leq\kappa^{-1}\|(1-S)QS W\|_{\mathrm{HS}}^2 .
\end{equation}
The factor $\kappa^{-1}$ is sharp. This gap concerns the norm
window in the unwindowed theta-image metric; it is not the
threshold used in the matrix concentration selection.
\end{proposition}
\begin{proof}
For $C=(1-S)QS$, the projection identities give
$C^*C=T-T^2$. On an eigenvector of $T$ with eigenvalue
$\lambda>0$, the first operator in
\eqref{glnw:inverse-boundary-criterion} has eigenvalue
$\lambda^{-1}\lambda(1-\lambda)=1-\lambda$; both sides
vanish on $\ker T$. Taking the finite-rank weighted trace gives
the Hilbert--Schmidt identity. If $\lambda\geq\kappa$, then
$1-\lambda\leq\kappa^{-1}\lambda(1-\lambda)$, proving
\eqref{glnw:window-gap-bound}. The rank-one construction below
attains equality when its window eigenvalue is $\kappa$.
Equivalently, if $Z$ is an isometry onto $\operatorname{Ran}Q$,
the gap is $Z^*SZ\geq\kappa\mathbf1_{(0,\infty)}(Z^*SZ)$.
\end{proof}

\begin{proposition}[Exterior smoothing does not determine normalization]
\label{glnw:normalization-insufficiency}
There are rank-one projections $Q_j$ and windows $S_j=S_{L_j}$,
$L_j\to\infty$, such that, for every $g,h\in C_c^\infty(\mathbb R)$,
\begin{equation}\label{glnw:left-smoothing-counterexample}
 \|V(g)(1-S_j)Q_j\|_{\mathrm{HS}}\to0,\qquad
 \Tr\bigl((Q_j-S_jQ_jS_j)V(h)\bigr)\to0,
\end{equation}
but $\Tr((\widetilde Q_j-S_jQ_jS_j)V(a)V(a)^*)$ stays positive
for a fixed $a\in C_c^\infty$. This is possible with the window
eigenvalue identically $1/2$. Without a uniform window gap, even
$\|(1-S_j)Q_jS_jV(g)\|_{\mathrm{HS}}\to0$ does not suffice.
If the fixed Godement--Jacquet data admit a nonzero norm-channel
seed, both obstructions also occur for sequences of theta
columns in the unrestricted matrix Schwartz image, preserving
all full Mellin zero jets of that seed. In the first theta example,
the window eigenvalue tends to $1/2$ and stays bounded away from zero.
No assertion is made that
these sequences belong to the prescribed matrix concentration
sections at $L=\log\Lambda$.
\end{proposition}
\begin{proof}
Choose unit vectors $u,v\in C_c^\infty((-1,1))$, put
$R_j=L_j+3$, and let $M_j\to\infty$. Define
\[
 v_j(t)=e^{-iM_jt}v(t-R_j),\qquad
 w_j=\sqrt c\,u+\sqrt{1-c}\,v_j,\qquad 0<c<1,
\]
and let $Q_j$ project onto $\mathbb Cw_j$. The summands are
orthogonal, $S_ju=u$, and $S_jv_j=0$. If $P_u$ projects onto
$\mathbb Cu$, direct calculation gives
\begin{equation}\label{glnw:rank-one-normalization}
 S_jQ_jS_j=cP_u,\qquad \widetilde Q_j=P_u,\qquad
 \widetilde Q_j-S_jQ_jS_j=(1-c)P_u.
\end{equation}
For every fixed $g\in C_c^\infty$, Plancherel gives
\[
 \|V(g)v_j\|_2^2
 =\frac1{2\pi}\int_{\mathbb R}
       |\widehat g(M_j+\eta)|^2|\widehat v(\eta)|^2\,d\eta\to0.
\]
Dominated convergence suffices; Schwartz decay gives arbitrarily
high inverse powers of $M_j$. The rank-one operator identities yield
\begin{align}
 \|V(g)(1-S_j)Q_j\|_{\mathrm{HS}}^2
   &=(1-c)\|V(g)v_j\|_2^2,\label{glnw:smoothing-order-left}\\
 \|(1-S_j)Q_jS_jV(g)\|_{\mathrm{HS}}^2
   &=c(1-c)\|V(g)^*u\|_2^2.\label{glnw:smoothing-order-right}
\end{align}
Also $\|(1-S_j)Q_jV(g)\|_{\mathrm{HS}}^2$ tends to the
right side of \eqref{glnw:smoothing-order-right}. Thus the
order of smoothing cannot be interchanged. For fixed compactly
supported $h$, the two cross terms in $\Tr(Q_jV(h))$ vanish
for large $j$ by support separation; its exterior diagonal term
tends to zero by the same Fourier estimate. This proves
\eqref{glnw:left-smoothing-counterexample}. In contrast,
\[
 \Tr((\widetilde Q_j-S_jQ_jS_j)V(a)V(a)^*)
       =(1-c)\|V(a)^*u\|_2^2,
\]
which is positive when $V(a)^*u\ne0$. Taking $c=1/2$ proves
the first assertion and taking $c=c_j\downarrow0$ proves the
failure without a gap in the right-ordered estimate. The ratio
of this normalization trace to the squared norm in
\eqref{glnw:smoothing-order-right}, with $g=a$, is exactly $c^{-1}$.

To obtain genuine theta columns, let $f_0=E_\pi\phi\ne0$ be a
seed from the fixed coefficient and finite data, with real matrix
input $\Phi$. Set $\omega_b=\omega_{\pi_\infty}(e^{b/n})$ and
\[
 \Phi^{[b]}(X)=e^{-nb/2}\omega_b^{-1}\Phi(e^{-b/n}X).
\]
This preserves the matrix Schwartz space and leaves the finite
data unchanged. Centrality and the substitution
$g=e^{b/n}g'$ in the Godement--Jacquet integral give
\[
 Z(s,\Phi^{[b]},\beta)
 =e^{-nb/2}\omega_b^{-1}\omega_b
       e^{b(s+(n-1)/2)}Z(s,\Phi,\beta)
 =e^{b(s-1/2)}Z(s,\Phi,\beta).
\]
The Mellin identity \eqref{glnw:mellin} and Fourier uniqueness
therefore give $E_\pi\Phi^{[b]}(t)=f_0(t-b)$, with the
pushforward understood. Rapid decay in both norm directions makes
$\widehat f_0$ entire and nonzero almost everywhere on $\mathbb R$.
The translates of $f_0$ span a dense subspace of $L^2(\mathbb R)$:
if $r$ is orthogonal to all of them, the $L^1$ function
$\widehat r\,\overline{\widehat f_0}$ has identically zero Fourier
transform, so it vanishes almost everywhere and hence $r=0$.

Choose finite linear combinations $f_j$ of these translates with
$\|f_j-w_j\|_2\leq\varepsilon_j=o(\sqrt{c_j})$, using constant
$c_j=c$ in the first example. Each $f_j$ is a theta
column with the same fixed data. After normalization, its
rank-one projection differs from $Q_j$ in trace norm by
$O(\varepsilon_j)$; since $\|S_jw_j\|_2=\sqrt{c_j}$, its
windowed support projection differs from $P_u$ by
$O(\varepsilon_j/\sqrt{c_j})$. These bounds follow by normalizing
the vectors and expanding the difference of their rank-one
projections. They preserve all the asserted limits for bounded
convolutions. Moreover the Mellin transform of $f_j$ is
\[
 Z_0(s)\sum_\ell c_{j,\ell}e^{b_{j,\ell}(s-1/2)},
\]
so every full zero jet of the seed remains exact. This density
argument bounds neither the required matrix Schwartz seminorms
nor their relation to a concentration scale; it proves no
membership assertion for the prescribed finite sections.
\end{proof}

The finite-rank positivity and quotient realization above are exact,
but they do not identify an all-zero harmonic-measure limit.
For $W=V(a)\otimes I_d$, the necessary normalization condition for the
selected sections is the vanishing of the norm in
\eqref{glnw:inverse-boundary-criterion}. A window gap and the
right-ordered estimate in \eqref{glnw:window-gap-bound} suffice;
left-smoothed exterior estimates alone do not. In fact the
windowed defect minus the unwindowed defect equals
\[
 \Tr((Q-SQS)WW^*)-\Tr((\widetilde Q-SQS)WW^*),
\]
so an outside-window trace comparison leaves this second term
undetermined. The stronger unsmoothed condition
$\|(1-S)Q\|_{\mathrm{HS}}\to0$ would also suffice by
\eqref{glnw:normalization-loss}, but is not asserted.
Windowing replaces exact Mellin jets by omitted-tail integrals.
Uniform control of these tails in the windowed Gram norm,
two boundary-closure identifications, local graph approximation,
and frequency-mass bounds remain necessary parts of a zero-measure
argument for the selected sections. Neither matrix concentration
nor the isometric quotient realization supplies these estimates
by itself.

